\section{Native induced residue frames and actual tensor powers}

\begin{lemma}\label{lem:actual-pulled-power-frame-naturality}
\lean{EnvelopeProperty.Berkovich.SchemeInducedPowerFrames.pullbackPowerFrame_naturality}
\leanok
\uses{lem:actual-power-pullback-naturality,lem:fc-actual-power-frame-iso-naturality,def:fc-actual-pulled-power-frame}
For every actual scheme map and original module isomorphism, the
genuine pulled power frame of the transported original frame is
the actual pulled power isomorphism followed by the original
pulled power frame, in every nonnegative degree.
\end{lemma}
\begin{proof}
Expand the pulled power frame as the native power pullback
comparison followed by the actual whole power frame. Naturality
of the native comparison moves the pulled power isomorphism past
that comparison. The whole power-frame naturality identity then
combines the original isomorphism with the original frame. The
result is exactly the pulled power frame of the transported
original frame. Every equality is a whole original isomorphism.
\end{proof}

\begin{lemma}\label{lem:actual-factor-power-frame-composition}
\lean{EnvelopeProperty.Berkovich.SchemeInducedPowerFrames.factorPowerFrame_comp}
\leanok
\uses{thm:fc-actual-pulled-power-frame-factor-coherence,thm:actual-power-pullback-composition,lem:fc-actual-power-frame-iso-naturality}
For an actual factorization $qb=tg$ and actual frame of $b^*M$,
the native composition comparison for $tg$ followed by the factor
transport of its pulled power frame equals the pullback of the
power comparison for $g$, followed by the pulled power frame for
$t$ of the actual transported degree-one frame.
\end{lemma}
\begin{proof}
The actual factor power identity first replaces factor transport
of the power frame by the power frame of the factor-transported
degree-one frame. Expand that pulled power frame. Native
composition coherence replaces the comparison for $tg$ by the
two successive power pullback comparisons and the power of the
native degree-one composition comparison. Whole power-frame
naturality combines this last power with the degree-one frame.
This is precisely the claimed entire frame identity.
\end{proof}

\begin{theorem}\label{thm:actual-induced-residue-power-frame}
\lean{EnvelopeProperty.Berkovich.SchemeInducedPowerFrames.inducedResidueFrame_postcomp,EnvelopeProperty.Berkovich.SchemeInducedPowerFrames.inducedResidueFrame_pullbackPower,EnvelopeProperty.Berkovich.SchemeInducedPowerFrames.inducedResidueFrame_power}
\leanok
\uses{lem:actual-factor-power-frame-composition,lem:actual-pulled-power-frame-naturality,def:fc-induced-native-residue-frame,lem:actual-power-iso-functoriality}
Let $g:Y\to Z$ be any actual scheme map, $y$ an actual native
valued point of $Y$, $\alpha:g^*M\simeq N$ an actual whole module
isomorphism, and $e$ an actual residue frame of $M$ at $g(y)$.
Raise $e$ through the genuine native power comparison, and
transport it through the canonical power isomorphism
$g^*(M^{\otimes n})\simeq N^{\otimes n}$ derived from $\alpha$.
The resulting induced native residue frame is exactly the genuine
pulled power frame of the induced degree-one residue frame.
This holds for all $n\geq0$ without properness, boundedness or a
frame-preservation premise.
\end{theorem}
\begin{proof}
Induced frames are the inverse actual pulled module isomorphism
followed by the genuine residue-module comparison and pullback
of the original frame along the actual residue-field map. Thus
postcomposing a module isomorphism postcomposes its induced frame
by the inverse pulled isomorphism. Apply the preceding factor
power lemma to the native residue-field square
$\operatorname{Spec}\kappa(y)\to\operatorname{Spec}\kappa(g(y))$.
The actual residue-module comparison is the native composition
comparison followed by this factor comparison. Cancel the actual
pulled power comparison for $g$. This proves the assertion when
the original isomorphism is the identity. For general $\alpha$,
the postcomposition identity separates its actual power from
that canonical comparison. Pulled power-frame naturality for
$\alpha^{-1}$ moves its inverse power into the original
degree-one frame; functoriality identifies this inverse with the
power of the inverse. The induced degree-one frame is exactly
this composite by its definition, proving the stated equality.
\end{proof}

\begin{corollary}\label{cor:actual-induced-normalized-power-frame}
\lean{EnvelopeProperty.Berkovich.SchemeInducedPowerFrames.inducedResidueFrame_normalizedPower}
\leanok
\uses{thm:actual-induced-residue-power-frame,lem:actual-power-iso-functoriality}
If $N=f^*L$ for any actual $f:Y\to X$, transport instead through
the normalized original power isomorphism to $f^*(L^{\otimes n})$.
Its induced native frame is the actual residue pullback of the
power comparison for $f$, followed by the genuine pulled power
frame of the same induced degree-one frame.
\end{corollary}
\begin{proof}
The normalized original power isomorphism is the preceding one
followed by the inverse genuine power pullback comparison for
$f$. The induced-frame postcomposition identity turns this
inverse into the direct residue pullback comparison. The
preceding theorem identifies the remaining entire frame with
the genuine power frame of the same degree-one induced frame.
\end{proof}
